\chapter{同余子群：怎样从有限余数回到精确数值}\label{ch:congruence}
\section{先看一个可以算到最后的输送}
给定两个本原整数列
\[v=\binom23,qquad w=\binom78.
\]
它们模五相同，但实际相差很大。矩阵
\[\gamma=\begin{pmatrix}-4&5\\-5&6\end{pmatrix}
\]
满足 $\det\gamma=1$、$\gamma\equiv I\pmod5$，并且 $\gamma v=w$。

这正是我们想要的工具：不改变指定余数，却能把实际向量精确送到另一个。定义
\[\Gamma(q,R)=\{M\in\SL_2(R):M\equiv I\pmod{qR}\}.
\]
$\SL_2(R)$ 是行列式为一的二阶矩阵群；“主同余子群”就是其中模 $q$ 等于单位矩阵的那些。这个名称虽然长，条件只是四个矩阵项的同余。

若 $R=\Z[1/D]$ 且 $\gcd(q,D)=1$，环里的分数怎样取模？把 $a/D^k$ 送到 $a(D^k)^{-1}\pmod q$。所以 $R/qR$ 与普通的 $\Z/q\Z$ 是同一个有限环。若 $q$ 含已经在 $R$ 中可逆的因子，先把它们去掉即可。例如在 $\Z[1/2]$ 中，模六的条件等价于模三，因为二已经是单位。

\section{为什么必须先取得这个群的真实实现}
模 $q$ 相同的两列有无穷多个。如果它们都能由 $\Gamma(q,R)$ 互相送到，而这个群的每个矩阵又都有真实平均词，那么无限问题就只剩有限问题：先送对余数，剩下由群精确完成。

不过，光在计算机里看到模 $q$ 的矩阵很多，并不能推出这个群已经可实现。模运算会把一大批不同有理矩阵压到一起，丢掉的部分正是我们需要控制的。这是两个外部定理真正进入证明的地方。

“有限指数”表示一个子群把整个群分成有限多个陪集，可以先把它理解成只剩有限种余类。Morris 的定理从足够多的初等矩阵得出有限指数；Serre 的定理再说，对我们的环，这样的子群一定包含某个主同余子群。
\begin{theorem}[Morris；Serre]\label{thm:arithmetic-external}
设 $D>1$、$R=R_D$。对于 $\SL_2(R)$ 的任一有限指数子群 $H$：
\begin{enumerate}
\item $H$ 中的上下初等矩阵生成 $H$ 的一个有限指数子群；
\item $H$ 包含某个非零层的主同余子群 $\Gamma(q,R)$。
\end{enumerate}
\end{theorem}
第一项使用 Morris\cite[Theorem 6.1(2)]{morris} 的二维无限单位情形。第二项使用 Serre\cite{serre} 的有理 $S$ 整数结论，准确陈述见\cite[pp.308--309]{prasad}。这里 $S$ 含实处和 $D$ 的素因子处，$D>1$ 保证至少有一个有限处。若 $D=1$，变成 $\SL_2(\Z)$，第二项一般就不成立。

这两项是本讲引用的深结果。后面从“不知道多大的模数”改进到明确的平方模数，则是我们自己能逐步算出的消元。

\section{为什么最后出现的是平方模数}
我们已经有上、下两族步长在 $hR$ 内的剪切。外部定理先保证一个未知 $\Gamma(q_0,R)$ 已包含其中。接下来想证明：任何 $g\equiv I\pmod{h^2}$ 也属于已生成的群。

取更深的模数 $q=h^2q_0$。在有限环 $R/qR$ 里，先给第一行加第二行的一个合适倍数，使左上角变成单位。然后像高斯消元一样消掉非对角项。剩下的对角项也能写成四个剪切，因为它与一的差被 $h^2$ 整除。

消元只保证与 $g$ 模 $q$ 相同。但剩余误差属于一开始已经取得的 $\Gamma(q_0,R)$，所以能够精确补掉。这就是“先有一个真群，再用有限运算改善它”的顺序。
\begin{lemma}[两条根生成一个确定的同余层]\label{lem:deep}
若 $D>1$、$h\ne0$，则
\[
\Gamma(h^2,R_D)\subseteq\langle U(hR_D),L(hR_D)\rangle.
\]
\end{lemma}
\begin{proof}

\par\noindent\textbf{先得到一个未知深层。}
记右侧为 $H$。将第一个外部定理用于有限指数子群 $\Gamma(h,R)$：
它包含的上下初等矩阵恰是参数在 $hR$ 中的剪切，故 $H$ 有限指数。
定理\ref{thm:arithmetic-external}先给出 $H$ 有限指数，再给出 $\Gamma(q_0,R_D)\subseteq H$。令 $q=h^2q_0$、$B=R_D/qR_D$。

\par\noindent\textbf{在有限环中找单位主元。}
对 $g=\left(\begin{smallmatrix}a&b\\c&d\end{smallmatrix}\right)\in\Gamma(h^2,R_D)$ 的模 $q$ 像，用 CRT 选 $k\in B$，使 $u=a+hkc$ 为单位：剩余特征整除 $h$ 时 $a=1$；其余分量若 $c\ne0$ 只排除一个 $k$ 类，若 $c=0$ 则 $a$ 已为单位。置 $v=b+hkd$，有
\[
g=U(-hk)L(c/u)\diag(u,u^{-1})U(v/u).
\]

从 $g\equiv I\pmod{h^2}$ 可知 $c\in h^2B$、$b\in h^2B$，
而 $v=b+hkd\in hB$。于是上式两侧的非对角消元都已经使用 $hB$ 中的参数。
还需拆开对角矩阵。由于 $u-1\in h^2B$，取 $t\in hB$ 使 $ht=u-1$。恒等式
\[
\diag(u,u^{-1})=U(h)L(t)U(-h/u)L(-tu)
\]
把中间项也分解为参数属于 $hB$ 的根。
\par\noindent\textbf{回到真实参数。}
约化 $hR_D\to hB$ 满射，因而可把有限个剪切参数逐个提升。
提升各参数至 $hR_D$，得到 $e\in H$ 与 $g$ 模 $q$ 相同。于是 $e^{-1}g\in\Gamma(q,R_D)\subseteq H$。
\end{proof}

\section{如果起初只有上剪切怎么办}
上剪切始终固定第一坐标轴。如果一个可逆返回 $J$ 把这条轴送到另一条直线，那么 $JU(t)J^{-1}$ 就固定那条新的直线。把两条直线当成坐标轴，就看到了上、下两种方向。

只是这个换基矩阵未必有整数逆，所以要把分母损失写进同余模数。下面的 $g^3$ 就是为这件事付出的明确代价。它不是多出三个物理位置，也不是三次平均的次数。
\begin{lemma}[横向返回与换基的分母损失]\label{lem:transverse}
假设 $U(hR)$ 以及 $J,J^{-1}$ 可实现，其中 $J=\left(\begin{smallmatrix}a&b\\c&d\end{smallmatrix}\right)\in M_2(R)$、$c\det J\ne0$。写 $c=g\epsilon$、$\epsilon\in R^\times$、$g\in\Z\setminus\{0\}$。则 $\Gamma(g^3h^2,R)$ 可实现。若有多个这样的 $g_i$，可将其层合并为 $\gcd(g_i)^3h^2$。
\end{lemma}
\begin{proof}
令 $P=(e_1,Je_1)=\left(\begin{smallmatrix}1&a\\0&c\end{smallmatrix}\right)$。相乘得
\[
P^{-1}U(t)P=U(ct),\qquad
P^{-1}JU(t)J^{-1}P=L(-ct/\det J).
\]
取 $t=hs$ 和 $t=\det J\,hs$，在 $P$ 基下得到两条完整 $ghR$ 根。引理\ref{lem:deep}给 $P\Gamma(g^2h^2,R)P^{-1}$。
因 $P,gP^{-1}$ 整于 $R$，若 $\gamma=I+g^3h^2X$，$X\in M_2(R)$，则
\[
P^{-1}\gamma P-I
=g^2h^2\,(gP^{-1})XP
\]
各项都属于 $g^2h^2R$，并且行列式仍为 $1$。
故 $P^{-1}\gamma P\in\Gamma(g^2h^2,R)$，证明该共轭群包含 $\Gamma(g^3h^2,R)$。
这个换基只用于论证群的包含关系，不需要把 $P$ 本身实现为平均词。
合并时先约化到各层的公倍数。在每个素数幂分量上，主同余条件嵌套，两个像生成较浅的条件；CRT 及初等矩阵提升给出各分量的乘积。接上已知公倍数层的核，即得最大公因子层。
\end{proof}

\section{相同余数的本原列，为什么真的能精确互达}
在 $R$ 中，“列 $(a,b)$ 本原”表示可以找到 $u,v\in R$ 使 $ua+vb=1$。在整数环这就是 $\gcd(a,b)=1$。它的线性代数含义是：这列能够补成一个行列式为一的二阶矩阵。

给起点 $v$ 补一列 $x$，给目标 $w$ 补一列 $y$。若能让两个补成的矩阵也模 $q$ 相同，那么
\[(w,y)(v,x)^{-1}
\]
便既把 $v$ 送到 $w$，又模 $q$ 等于 $I$。下面证明的全部工作，就是把第二列稍作修改，以满足这两个要求。
\begin{example}[同余相同，实际数值可以相差很大]
取 $v=(1,0)^T$、$w=(1,q)^T$，两者模 $q$ 相同。
下剪切 $L(q)$ 已经把 $v$ 精确送到 $w$。
一般本原列可能没有这么直接的一个剪切，但下一个引理保证某个 $\Gamma(q,R)$ 元素总能完成。
\end{example}
\begin{lemma}[相同剩余类中的精确输送]\label{lem:fiber}
若 $v,w\in R^2$ 本原且 $v\equiv w\pmod q$，则存在 $\gamma\in\Gamma(q,R)$ 使 $\gamma v=w$。
\end{lemma}
\begin{proof}
取 $x$ 使 $\det(v,x)=1$，写 $\det(w,x)=1+qk$。
写 $w=(w_1,w_2)^T$，则 $\det(w,z)=w_1z_2-w_2z_1$。
因为 $w_1,w_2$ 生成整个 $R$，这个线性式能够取到任意环元素。
也就是说，本原性使 $z\mapsto\det(w,z)$ 满射，故取 $z$ 使其值为 $-k$。令 $y=x+qz$，则 $(w,y)\equiv(v,x)\pmod q$，两矩阵行列式皆为 $1$。取 $\gamma=(w,y)(v,x)^{-1}$。
\end{proof}
这是整套方法里最值得停下来理解的一步：有限同余计算负责把你送进正确的“房间”，这个引理负责在房间内送到那个\emph{确切的点}。最终是等号，不是取极限。

\section{取得主同余子群以后，逆更容易造了}
在有限矩阵群里，某矩阵的正整数次幂一定回到 $I$。如果 $Z^m\equiv I\pmod q$，那么真实 $Z^m$ 就落进已能实现的 $\Gamma(q,R)$，其逆也已可用。因此
\[Z^{-1}=Z^{m-1}(Z^m)^{-1}.
\]
这把“有限商中的周期”变成了一条真正的逆词。

行列式不是一时，要额外处理整数内容。Smith 标准形是通过整数可逆行列变换把矩阵化成对角形的定理；这里用它设计一个两边夹乘，使结果能除掉行列式而仍是整数矩阵。再由 CRT 让这个夹乘同时模 $q$ 正确。
\begin{lemma}[获得深同余群以后的正向逆]\label{lem:inverse}
若 $\Gamma(q,R)$ 可实现，则每个可实现 $Z\in\SL_2(R)$ 的逆可实现。若整数满秩返回 $A$ 满足 $\gcd(\det A,q)=1$，则 $A^{-1}$ 亦有射影实现。
\end{lemma}
\begin{proof}
第一项取 $m$ 使 $Z^m\equiv I\pmod q$，用 $Z^{-1}=Z^{m-1}(Z^m)^{-1}$。
第二项取偶数 $m$ 使 $P=A^m\equiv I\pmod q$，令 $\Delta=\det P>0$。选 $L,R'\in\SL_2(\Z)$ 使 $LPR'=\diag(a,b)$、$ab=\Delta$。置 $J_0=\left(\begin{smallmatrix}0&1\\-1&0\end{smallmatrix}\right)$、$g_*=R'J_0L$，则 $Pg_*P$ 各项被 $\Delta$ 整除。
利用 CRT 及 $\SL_2(\Z)$ 对有限模数的满射，取 $g\equiv g_*\pmod\Delta$、$g\equiv I\pmod q$。则 $H=\Delta^{-1}PgP\in\Gamma(q,\Z)$，而 $gPH^{-1}=\Delta P^{-1}$。接上 $A^{m-1}$ 得结论。所用约化满射由有限局部环上的初等消元、CRT 和整数参数提升得到。
\end{proof}
第一次读到这儿，可先记住它的用途和条件：\textbf{已经取得真实主同余子群后，行列式与控制模数互素的整数返回，可以获得正向射影逆。}具体夹乘等式留在上面的完整证明中，第二遍逐项相乘核对即可。
